\section{Death and Life in Wondrous Duel: Christ Enters the Grave}

On Easter morning, before the Gospel of the empty tomb is proclaimed, the Roman liturgy breaks into a blazing sequence that states the turning point of human history in four Latin lines:

\begin{quote}
\emph{Mors et vita duello conflixere mirando: dux vitae mortuus regnat vivus.}\endnote{1962 \emph{Missale Romanum}, Dominica Resurrectionis, Sequence \emph{Victimae paschali laudes}, p.~330.}
\end{quote}

``Death and life contended in a wondrous duel: the Prince of life, who died, reigns alive.'' Every Christian word about mortality rests on that duel. God did not repeal the sentence of death by a distant decree, leaving man's debt unpaid and the grave an unentered terror. Instead, the eternal Son took flesh capable of dying, suffered the real separation of soul and body upon the Cross, descended to the righteous prisoners of the old world, and broke the lock from within.

\subsection{The Word Takes a Mortal Body to Pay the Common Debt}

The Epistle to the Hebrews places the conquest of death at the center of the Incarnation: ``Therefore because the children are partakers of flesh and blood, he also himself in like manner hath been partaker of the same: that, through death, he might destroy him who had the empire of death, that is to say, the devil: and might deliver them, who through the fear of death were all their lifetime subject to servitude'' (Heb.\ 2:14--15). The weapon is \emph{through death} (\emph{per mortem}): death is not bypassed, but entered and undone by its own stroke.

No Father expounded this mystery with greater vigor than St.~Athanasius of Alexandria in \emph{On the Incarnation of the Word}. Unable to bear that His Father's handiwork should perish under corruption, the immortal Word took a body ``of no different sort from ours'' from a spotless Virgin so that He might offer it as a sacrifice in our stead:

\begin{quote}
For the Word, perceiving that no otherwise could the corruption of men be undone save by death as a necessary condition, while it was impossible for the Word to suffer death, being immortal, and Son of the Father; to this end He takes to Himself a body capable of death, that it, by partaking of the Word Who is above all, might be worthy to die in the stead of all, and might, because of the Word which was come to dwell in it, remain incorruptible, and that thenceforth corruption might be stayed from all by the Grace of the Resurrection.\endnote{St.~Athanasius of Alexandria, \emph{On the Incarnation of the Word}, \S~9.1 (Archibald Robertson trans., NPNF, ser.~2, vol.~4, 1892), pp.~40--41; cf.~\S~8.2--3, \S~9.3--4, and \S~20.5 (``two marvels came to pass at once, that the death of all was accomplished in the Lord's body, and that death and corruption were wholly done away by reason of the Word that was united with it'').}
\end{quote}

Like a king who enters a besieged city and dwells in a single house, rescuing the whole city from bandits, the Monarch of all entered one human body and halted death's ancient conspiracy against the human race. The result was not an abstract theorem, but a visible revolution in human courage:

\begin{quote}
For as when a tyrant has been defeated by a real king, and bound hand and foot, then all that pass by laugh him to scorn, buffeting and reviling him, no longer fearing his fury and barbarity, because of the king who has conquered him; so also, death having been conquered and exposed by the Saviour on the Cross, and bound hand and foot, all they who are in Christ, as they pass by, trample on him, and witnessing to Christ scoff at death, jesting at him, and saying what has been written against him of old: ``O death, where is thy victory? O grave, where is thy sting?''\endnote{St.~Athanasius, \emph{On the Incarnation}, \S~27.4 (NPNF2 vol.~4, p.~51); cf.~\S\S~28--29 (comparing conquered death to a slain serpent or dead lion).}
\end{quote}

\subsection{The Truth of Christ's Death and the Unbroken Union}

To guard that victory from every dilution, Catholic faith holds two truths together concerning Good Friday and Holy Saturday: Christ truly and voluntarily died a real human death, yet His Divinity never parted from either His departed soul or His lifeless body.

First, His death was real and sovereignly voluntary. He did not yield to disease or old age as a victim overpowered by nature: ``No man taketh it away from me: but I lay it down of myself. And I have power to lay it down: and I have power to take it up again'' (John 10:18; cf.\ 10:17).\endnote{Aquinas, \emph{Summa Theologiae}, III, q.~50, a.~1, ad 2; St.~Athanasius, \emph{On the Incarnation}, \S~21.4.} In Gethsemane, He genuinely felt the shudder of human flesh---``My soul is sorrowful even unto death'' (Matt.\ 26:38)---and Pope St.~Leo the Great hails that trembling as a merciful exchange: ``In our Nature, therefore, the Lord trembled with our fear, that He might fully clothe our weakness and our frailty with the completeness of His own strength. For He had come into this world a rich and merciful Merchant from the skies, and by a wondrous exchange had entered into a bargain of salvation with us, receiving ours and giving His, honour for insults, salvation for pain, life for death.''\endnote{St.~Leo the Great, \emph{Sermon LIV} (On the Passion, III), \S\S~4--5 (Charles Lett Feltoe trans., NPNF, ser.~2, vol.~12, p.~166).} On Calvary, crying aloud in the words of Psalm 30---``Father, into thy hands I commend my spirit''---He bowed His head and gave up the ghost (Luke 23:46; John 19:30).

In the \emph{Tertia Pars} of the \emph{Summa Theologiae}, St.~Thomas Aquinas gives five reasons why it was fitting for Christ to die: to satisfy for a race sentenced to death; to prove the reality of His flesh; to free us from lifelong slavery to the fear of death; to set the pattern of dying spiritually to sin; and by rising from a real grave to kindle our hope of resurrection. Because a living man consists in the union of soul and body, during the three days in the tomb Christ's human nature was genuinely sundered, so that He was truly ``a dead man''---confessed in the Creed as ``crucified, dead, and buried.''\endnote{Aquinas, \emph{Summa Theologiae}, III, q.~50, aa.~1, 4; cf.\ \emph{Roman Catechism}, pt.~1, Creed Art.~IV (Donovan, p.~45).}

Yet death could never touch the hypostatic union. ``What is bestowed through God's grace is never withdrawn except through fault,'' Thomas observes; quoting St.~John of Damascus, he affirms that although Christ's holy soul was separated from His spotless body, ``His Godhead remained unseparated from both---from the soul, I mean, and from the body.''\endnote{Aquinas, \emph{Summa Theologiae}, III, q.~50, a.~2, and a.~3, \emph{sed contra} and ad 4; cf.\ Fourth Lateran Council (1215), Constitution 1, \emph{Firmiter credimus} (DS 801): \emph{descendit in anima, et resurrexit in carne} (He descended in soul, and rose in flesh).} St.~Leo the Great preaches the same mystery on Easter: ``the Godhead, Which quitted not either part of the Human Nature which He had assumed, reunited by Its power that which Its power had separated.''\endnote{St.~Leo the Great, \emph{Sermon LXXI} (On the Lord's Resurrection, I), \S~2 (NPNF2 vol.~12, p.~182).} The \emph{Roman Catechism} states this plainly for the faithful:

\begin{quote}
When, therefore, we say that Jesus died, we mean that his soul was disunited from his body; not that his divinity was so separated. On the contrary, we firmly believe and profess that, when his soul was dissociated from his body, his divinity continued always united both to his body in the sepulchre, and to his soul in Limbo. [\ldots] For, as the divinity was never separated from his body which was laid in the sepulchre, we truly confess that God was buried.\endnote{\emph{Roman Catechism}, pt.~1, Creed Art.~IV (Donovan, pp.~45--46).}
\end{quote}

United personally to the Word, Christ's dead body saw no decay, fulfilling Psalm 15:10 (``Nor wilt thou give thy holy one to see corruption''), and remained ``the instrument of the Godhead united to Him, working by Its power, although dead''---so that, in St.~Augustine's phrase quoted by Aquinas, the one bodily death of our Savior saved us from our two deaths of soul and body.\endnote{Aquinas, \emph{Summa Theologiae}, III, q.~50, aa.~5--6 and ad 3; \emph{Roman Catechism}, pt.~1, Creed Art.~IV (p.~46).}

\subsection{The Descent into the Limbo of the Fathers}

In the silence of Holy Saturday, Christ's human soul---united to His Godhead---descended ``into hell'' (\emph{descendit ad inferos}). As the \emph{Roman Catechism} explains under the fifth article of the Creed, this denotes not the grave already confessed, nor the torment of the damned, but the hidden abodes of souls not yet admitted to the beatific vision---specifically the Limbo of the Fathers (\emph{Limbus Patrum}), ``Abraham's bosom'' (Luke 16:22), where the righteous of the Old Covenant rested without sensible pain, awaiting the price of redemption. ``They descended as captives,'' says the Catechism, ``He as free and victorious amongst the dead, to subdue those demons by whom, in consequence of primeval guilt, they were held in captivity.''\endnote{\emph{Roman Catechism}, pt.~1, Creed Art.~V (Donovan, pp.~50--52); cf.\ Aquinas, \emph{Summa Theologiae}, III, q.~52, aa.~1--2.}

St.~Thomas shows how that descent fulfilled the ancient prophecies: Osee's challenge, ``O death, I will be thy death; O hell, I will be thy bite'' (Os.\ 13:14), and Zacharias's promise, ``Thou also by the blood of thy testament hast sent forth thy prisoners out of the pit, wherein is no water'' (Zach.\ 9:11). Entering the strong man's house, Christ poured the light of eternal glory upon Adam, Abel, Abraham, Moses, David, the prophets, and the penitent thief promised ``This day thou shalt be with me in paradise'' (Luke 23:43), shattering the bars of the ancient prison so that every knee in heaven, on earth, and under the earth might bow to His name.\endnote{Aquinas, \emph{Summa Theologiae}, III, q.~52, aa.~1, 4 (ad 3), and 5; \emph{Roman Catechism}, pt.~1, Creed Art.~V (p.~52).}

When the third day dawned, Christ rose as ``the firstfruits of them that sleep,'' for ``as in Adam all die, so also in Christ all shall be made alive'' (1 Cor.\ 15:20, 22). To St.~John on Patmos, the Risen Lord proclaims His conquered title: ``I am living for ever and ever and have the keys of death and of hell'' (Apoc.\ 1:18). Though ``the enemy, death, shall be destroyed last'' (1 Cor.\ 15:26), its empire is broken, and the Apostle chants the triumph of the redeemed:

\begin{quote}
And when this mortal hath put on immortality, then shall come to pass the saying that is written: Death is swallowed up in victory. O death, where is thy victory? O death, where is thy sting? Now the sting of death is sin: and the power of sin is the law. But thanks be to God, who hath given us the victory through our Lord Jesus Christ. (1 Cor.\ 15:54--57)
\end{quote}

Because sin was death's sting, and Christ has drawn that sting into His sinless flesh upon the Cross, bodily death no longer meets the believer as an undefeated master.
